% 可复核的二次函数：共同初值 w=(1,2), w'=(3,2)，步长 eta=1/2。
% 三组 Hessian 分别为 diag(-1,1), diag(0,1), diag(1,2)。
% 更新后的点分别为 (1.5,1),(4.5,1)；(1,1),(3,1)；(.5,0),(1.5,0)。
% 三个参数平面共用 8.5 mm/单位的尺度；初始距离均为 2。
\begin{tikzpicture}[foundation picture]
 \path[use as bounding box] (0,-1) rectangle (169,70);
 \foreach \offset/\casecolor/\title/\formula/\result/\condition in {
  0/FigureBlue/仅光滑/{f=-\tfrac12w_1^2+\tfrac12w_2^2}/{D^+=3=1.5D}/{M=1},
  57/FigureRed/凸且光滑/{f=\tfrac12w_2^2}/{D^+=2=D}/{M=1,\ \eta\leq2/M},
  114/FigureYellow/强凸且光滑/{f=\tfrac12w_1^2+w_2^2}/{D^+=1=0.5D}/{\mu=1,\ M=2,\ \eta\leq1/M}}{
  \begin{scope}[shift={(\offset,0)}]
   \node[foundation heading] at (25,66) {\title};
   \node at (25,57) {$\formula$};
   \draw[foundation axis] (5,26)--(48,26) node[right] {$w_1$};
   \draw[foundation axis] (5,26)--(5,50) node[above] {$w_2$};
   \node[below left] at (5,26) {$0$};
   \draw[FigureMuted,line width=.6pt,dashed] (13.5,43)--(30.5,43);
   \node at (22,49) {$D=2$};
   \node[anchor=south east] at (12,44) {$\bm w$};
   \node[anchor=south west] at (32,44) {$\bm w'$};
   \node[text=FigureInk] at (25,12) {$\result$};
   \node at (25,4) {$\condition$};
  \end{scope}
 }
 % 光滑非凸：空心/实心圆标记更新前后。
 \draw[FigureBlue,line width=1.1pt,-{Stealth[length=1.8mm]}]
  (13.5,43)--(17.75,34.5);
 \draw[FigureBlue,line width=1.1pt,-{Stealth[length=1.8mm]}]
  (30.5,43)--(43.25,34.5);
 \draw[FigureBlue,line width=1.1pt] (17.75,34.5)--(43.25,34.5);
 \foreach \x in {13.5,30.5}{\draw[FigureBlue,fill=white,line width=1.1pt] (\x,43) circle (1.2);}
 \foreach \x in {17.75,43.25}{\fill[FigureBlue] (\x,34.5) circle (1.2);}
 \node[below] at (17.75,33) {$\bm w^+$};
 \node[below] at (43.25,33) {$\bm w'^{+}$};
 % 凸：方形标记；平坦的 w_1 方向使本例的距离保持不变。
 \begin{scope}[shift={(57,0)}]
  \foreach \x in {13.5,30.5}{
   \draw[FigureRed,line width=1.1pt,-{Stealth[length=1.8mm]}] (\x,43)--(\x,34.5);
   \draw[FigureRed,fill=white,line width=1.1pt] (\x-1.1,41.9) rectangle (\x+1.1,44.1);
   \fill[FigureRed] (\x-1.1,33.4) rectangle (\x+1.1,35.6);
  }
  \draw[FigureRed,line width=1.1pt] (13.5,34.5)--(30.5,34.5);
  \node[below] at (13.5,33) {$\bm w^+$};
  \node[below] at (30.5,33) {$\bm w'^{+}$};
 \end{scope}
 % 强凸：三角形标记；实线终点距离为初始虚线的一半。
 \begin{scope}[shift={(114,0)}]
  \draw[FigureYellow,line width=1.1pt,-{Stealth[length=1.8mm]}] (13.5,43)--(9.25,26);
  \draw[FigureYellow,line width=1.1pt,-{Stealth[length=1.8mm]}] (30.5,43)--(17.75,26);
  \draw[FigureYellow,line width=1.1pt] (9.25,26)--(17.75,26);
  \foreach \x in {13.5,30.5}{
   \draw[FigureYellow,fill=white,line width=1.1pt] (\x,44.4)--(\x-1.3,42.3)--(\x+1.3,42.3)--cycle;
  }
  \foreach \x in {9.25,17.75}{\fill[FigureYellow] (\x,27.4)--(\x-1.3,25.3)--(\x+1.3,25.3)--cycle;}
  \node[below] at (8,24.4) {$\bm w^+$};
  \node[below] at (20,24.4) {$\bm w'^{+}$};
 \end{scope}
\end{tikzpicture}%
